\section*{अध्याय \ref{ch:b2:isotope-geochemistry} --- समस्थानिक भू-रसायन}

\begin{solution}{exo:b2:isotope-geochemistry:1}
ट्रिस्टन: $(0.51249/0.512638 - 1)\times10^4 = -2.9$। मंगाइया: $(0.51285/0.512638 -
1)\times10^4 = +4.1$। MORB: $+8.5$। मंगाइया और MORB के स्रोतों का, जिनका
$\varepsilon_{\mathrm{Nd}}$ धनात्मक है, Sm/Nd औसतन \omterm{def:b2:planet-formation-meteorites:chondrule}{कॉन्ड्राइटों} से ऊँचा रहा है;
ट्रिस्टन के स्रोत का निम्न।
\end{solution}

\begin{solution}{exo:b2:isotope-geochemistry:2}
$\delta^{18}\mathrm{O}_{\mathrm{VPDB}} = (\num{2.0672e-3}/\num{2.0052e-3} - 1)\times1000 =
+30.9$ प्रति मिल, \omterm{def:b2:isotope-geochemistry:standards}{VSMOW} पैमाने पर। कैल्साइट का $^{18}\mathrm{O}/^{16}\mathrm{O} =
\num{2.0672e-3}\times(1 - 0.0020) = \num{2.0631e-3}$ है, अतः $\delta^{18}\mathrm{O} =
(\num{2.0631e-3}/\num{2.0052e-3} - 1)\times1000 = +28.9$ प्रति मिल (\omterm{def:b2:isotope-geochemistry:standards}{VSMOW})।
\end{solution}

\begin{solution}{exo:b2:isotope-geochemistry:3}
\qty{1.0}{Ga} पर: $e^{\lambda t} - 1 = 0.006545$, और $0.512638 - 0.1966\times0.006545 =
0.511351$। \qty{4.567}{Ga} पर: $e^{\lambda t}-1 = 0.030243$, और $0.512638 -
0.1966\times0.030243 = 0.506692$।
\end{solution}

\begin{solution}{exo:b2:isotope-geochemistry:4}
$\delta\mathrm{D} = 8\times(-18.0) + 10 = -134$ प्रति मिल। झील के लिए, $d = -42 -
8\times(-4.0) = -10$ प्रति मिल, रेखा के 10 से नीचे: जल \omterm{prop:b2:isotope-geochemistry:mwl}{उल्कीय
जल रेखा} के नीचे है, जैसे वाष्पित पृष्ठीय जल होते हैं (हल्के समस्थानिक पहले निकल गए)।
\end{solution}

\begin{solution}{exo:b2:isotope-geochemistry:5}
$Q = 10^4\times\num{6.524e-12}\times0.1966/0.512638 = \num{2.502e-8}$ प्रति वर्ष, या
\qty{25.0}{\per\giga\year}। ग्रेनाइट का $\mathrm{Sm}/\mathrm{Nd} = 6.0/30 = 0.200$ है,
CI के $0.145/0.457 = 0.317$ के सामने: $f = 0.200/0.317 - 1 = -0.37$।
\qty{0.5}{Gyr} में: $\Delta\varepsilon = 25.0\times(-0.37)\times0.5 = -4.6$ इकाई।
\end{solution}

\begin{solution}{exo:b2:isotope-geochemistry:6}
$\varepsilon \approx 3.5 + 25.0\times(-0.45)\times2.7 = 3.5 - 30.4 = -26.9$,
$-27.3$ के सामने: सन्निकटन 0.4 इकाई अधिक ऊँचा है (पर्याप्त ऋणात्मक नहीं)। अंतर का
अधिकांश, लगभग 0.3 इकाई, $e^{\lambda t} > 1$ से आता है: जनक, और इसलिए
CHUR के सापेक्ष वृद्धि की कमी, अतीत में आज से अधिक प्रचुर थी
(\cref{exo:b2:isotope-geochemistry:10}); शेष $f$ और आरंभिक मान के पूर्णांकन से आता है।
\end{solution}

\begin{solution}{exo:b2:isotope-geochemistry:7}
ग्रेनाइट के द्रव्यमान-अंश $x$ के साथ, $\mathrm{Sr} = 320x + 130(1-x)$ और
$R = (320x\times0.7200 + 130(1-x)\times0.7030)/\mathrm{Sr}$। $x = 0.05$ के लिए: Sr
\qty{139.5}{ppm}, $R = 0.70495$; $x = 0.10$: 149.0, 0.70665; $x = 0.20$: 168.0,
0.70948। मिश्रण प्रतिज्ञप्ति की उपपत्ति से, $R\,\mathrm{Sr}$ Sr में रैखिक है:
$R = 0.731632 - 3.7221/\mathrm{Sr}$, प्रवणता $-3.722$ (ppm) और अंतःखंड 0.7316 वाली
एक सरल रेखा, $1/\mathrm{Sr}$ के सापेक्ष। $x = 0.05$ पर जाँच: $0.731632 - 3.7221/139.5
= 0.70495$।
\end{solution}

\begin{solution}{exo:b2:isotope-geochemistry:8}
\qty{20}{\degreeCelsius} पर: $\delta^{18}\mathrm{O} = 1000/1.0098 - 1000 = -9.70$ प्रति
मिल, $\delta\mathrm{D} = 1000/1.084 - 1000 = -77.5$ प्रति मिल, प्रवणता $77.5/9.70 = 7.99$।
\qty{0}{\degreeCelsius} पर: $-11.56$ और $-99.9$ प्रति मिल, प्रवणता 8.64। दोनों \omterm{prop:b2:isotope-geochemistry:mwl}{उल्कीय
जल रेखा} की प्रवणता 8 के निकट हैं: द्रव और वाष्प के बीच जल का साम्य प्रभाजन
अकेले ही ड्यूटेरियम और ऑक्सीजन-18 के विस्थापनों का अनुपात
8 के निकट रखता है।
\end{solution}

\begin{solution}{exo:b2:isotope-geochemistry:9}
भूपर्पटी $0.0051\times27/1.25 = 0.110$, अर्थात् नियोडिमियम का 11 प्रतिशत धारण करती है।
द्रव्यमान-संतुलन: $0.110\times(-15) + 0.890\,\varepsilon_m = 0$, अतः $\varepsilon_m = 1.65/0.890 =
+1.9$। हिस्सा $x$ धारण करने वाले, $+8.5$ पर स्थित, क्षीणित भंडार और शेष 0 पर होने से:
$0.110\times(-15) + 8.5x = 0$, $x = 1.65/8.5 = 0.19$: नियोडिमियम का लगभग
पाँचवाँ भाग।
\end{solution}

\begin{solution}{exo:b2:isotope-geochemistry:10}
वृद्धि-दर $\lambda p\,e^{\lambda t}$ है, जहाँ $e^{\lambda t} \approx 1 + \lambda t$;
अलगाव $T$ से वर्तमान तक के अंतराल पर औसत लेने पर गुणक
$1 + \lambda T/2$ है, अतः उपेक्षित पदों का सापेक्ष आकार $\lambda T/2$ है। पृथ्वी की
आयु में, $\num{6.524e-12}\times\num{4.567e9}/2 = 0.015$, अर्थात् 1.5 प्रतिशत।
\cref{exo:b2:isotope-geochemistry:6} की भूपर्पटी के लिए $\lambda T/2 = 0.0088$, और त्रुटि $0.0088\times30.4 = 0.27$
इकाई है। CHUR अनुपात का धीमा परिवर्तन, जिससे $\varepsilon$ में भाग दिया जाता है,
एक और छोटा पद जोड़ता है, $(p/R)\lambda T \times \varepsilon_0$ की कोटि का, यहाँ 0.02 इकाई।
\end{solution}

\begin{solution}{exo:b2:isotope-geochemistry:11}
$K = 13.3/5.0 = 2.66$। $s = 0.5$ के साथ: $s' = 2.66\times0.5/(1 + 1.66\times0.5) =
1.33/1.83 = 0.73$। अंतर $s' - s$ वहाँ सबसे बड़ा है जहाँ
$\mathrm{d}s'/\mathrm{d}s = K/(1+(K-1)s)^2 = 1$, अर्थात् $s = (\sqrt K - 1)/(K-1) =
0.631/1.66 = 0.38$ पर, जहाँ $s' = 0.62$: 0.24 का अंतर।
\end{solution}

\begin{solution}{exo:b2:isotope-geochemistry:12}
जनन: $2.3$ से $2.6\times\qty{7.2e9}{km^3}/\qty{4.5e9}{yr} = 3.7$ से
\qty{4.2}{km^3} प्रति वर्ष। विनाश: $1.3$ से $1.6\times7.2/4.5 = 2.1$ से
\qty{2.6}{km^3} प्रति वर्ष। वर्तमान दरें भिन्न हैं: जुड़ाव (2.7) मॉडल के
माध्य जनन से नीचे है, हटाव (3.2) उसके माध्य विनाश से ऊपर। मॉडल में भूपर्पटी
आरंभ में बढ़ी, जब जनन सबसे तेज़ था, और वर्तमान संतुलन या थोड़ी
हानि का समय है, जैसा क्षेपण-क्षेत्र का हिसाब बताता है।
\end{solution}

\begin{solution}{pb:b2:isotope-geochemistry:1}
\textbf{1.} प्रावार: $(0.51307/0.512638 - 1)\times10^4 = +8.4$। अवसाद:
$(0.51141/0.512638 - 1)\times10^4 = -24.0$।

\textbf{2.} कोरा $+2.8$, वेल $-4.3$, आशा $+4.7$।

\textbf{3.} $^{87}\mathrm{Sr}/^{86}\mathrm{Sr}$ क्षैतिज और $\varepsilon_{\mathrm{Nd}}$ ऊर्ध्वाधर लेकर:
प्रावार (0.70282, $+8.4$), कोरा (0.70410, $+2.8$), वेल (0.70618, $-4.3$), अवसाद
(0.72000, $-24.0$), आशा (0.70290, $+4.7$)। केवल वेल, ऋणात्मक
$\varepsilon_{\mathrm{Nd}}$ के साथ, CHUR के समृद्ध पक्ष पर है।

\textbf{4.} सरणी के नीचे: \omterm{def:b2:isotope-geochemistry:reservoirs}{क्षीणित प्रावार} के स्ट्रॉन्शियम अनुपात पर सरणी का
$\varepsilon_{\mathrm{Nd}}$ लगभग $+8$ होता; आशा का $+4.7$ है। यह
HIMU-प्रकार के सेंट हेलेना समूह (मंगाइया, तुबुआई, सेंट हेलेना) से मिलता-जुलता है।

\textbf{5.} गलन जनक--संतति अनुपात बदलता है पर स्ट्रॉन्शियम और नियोडिमियम के
समस्थानिक अनुपात मापने योग्य रूप से नहीं (गलन वाली प्रतिज्ञप्ति); लावा
युवा हैं, इसलिए उद्गार के बाद से उनके अनुपात नहीं बढ़े (प्रश्न 6)।

\textbf{6.} $\Delta R = 0.12\,(e^{\lambda\cdot10^6} - 1) = 0.12\times\num{6.52e-6} =
\num{7.8e-7}$, या $\num{7.8e-7}/0.512638\times10^4 = 0.015$ एप्सिलॉन इकाई: द्वीपों के बीच के
अंतरों से बहुत नीचे, और किसी मापन की परिशुद्धता से भी नीचे।

\textbf{7.} $\mathrm{Nd} = 40f + 0.75(1-f)$, और
$R = \bigl(40f\times0.51141 + 0.75(1-f)\times0.51307\bigr)/\mathrm{Nd}$।

\textbf{8.} $f = 0.5$ प्रतिशत: Nd \qty{0.9463}{ppm}, $\varepsilon = +1.6$; 1 प्रतिशत:
1.1425, $-2.9$; 2 प्रतिशत: 1.535, $-8.4$; 5 प्रतिशत: 2.7125, $-15.4$।

\textbf{9.} Sr 10.95, 11.90, 13.80 और \qty{19.50}{ppm}; $^{87}\mathrm{Sr}/^{86}\mathrm{Sr} =
0.70439$, 0.70571, 0.70780 और 0.71163।

\textbf{10.} $R\,\mathrm{Sr} = 10\times0.70282 + f(200\times0.72000 - 10\times0.70282)$ और
$\mathrm{Sr} = 10 + 190f$; $f$ का विलोपन करने पर, $R\,\mathrm{Sr} = 7.0282 +
(\mathrm{Sr}-10)\times136.9718/190$, अतः $R = 0.720904 - 0.18084/\mathrm{Sr}$: प्रवणता $-0.1808$ (ppm) वाली एक सरल
रेखा, $1/\mathrm{Sr}$ के सापेक्ष। $f = 1$ प्रतिशत पर जाँच:
$0.720904 - 0.18084/11.90 = 0.70571$।

\textbf{11.} $K = (10/0.75)/(200/40) = 13.33/5.0 = 2.67$। प्रावार सापेक्षतः
स्ट्रॉन्शियम में अधिक समृद्ध है, अवसाद नियोडिमियम में, इसलिए Nd अनुपात
अवसाद की ओर तेज़ी से चलता है: थोड़ा-सा अवसाद $\varepsilon_{\mathrm{Nd}}$ को बहुत दूर खिसका देता है, इससे पहले कि
$^{87}\mathrm{Sr}/^{86}\mathrm{Sr}$ अधिक चले।

\textbf{12.} Nd में 53 गुना समृद्ध ($40/0.75$), Sr में 20 गुना। $f = 1$ प्रतिशत पर
अवसाद नियोडिमियम का $0.40/1.1425 = 35$ प्रतिशत देता है, जो
$\varepsilon_{\mathrm{Nd}}$ को $0.35\times(8.4 + 24.0) = 11.3$ इकाई खिसकाता है, पर स्ट्रॉन्शियम का केवल $2.0/11.9 =
17$ प्रतिशत।

\textbf{13.} $0.01\times40/1.1425 = 0.350$: 35 प्रतिशत।

\textbf{14.} $R\,(40f + 0.75(1-f)) = 40fR_S + 0.75(1-f)R_D$ से:
$f = 0.75(R - R_D)/\bigl(40(R_S - R) + 0.75(R - R_D)\bigr)$, और स्ट्रॉन्शियम के लिए यही, 200
और 10 के साथ।

\textbf{15.} $f = 0.75\times(-0.00029)/\bigl(40\times(-0.00137) + 0.75\times(-0.00029)\bigr)
= \num{-2.175e-4}/\num{-5.502e-2} = 0.0040$: 0.40 प्रतिशत।

\textbf{16.} $f = 10\times0.00128/\bigl(200\times0.01590 + 10\times0.00128\bigr) =
0.0128/3.1928 = 0.0040$: 0.40 प्रतिशत; दोनों मेल खाते हैं।

\textbf{17.} नियोडिमियम: $0.75\times(-0.00065)/\bigl(40\times(-0.00101) +
0.75\times(-0.00065)\bigr) = \num{-4.875e-4}/\num{-4.089e-2} = 0.0119$; स्ट्रॉन्शियम:
$10\times0.00336/\bigl(200\times0.01382 + 10\times0.00336\bigr) = 0.0336/2.7976 = 0.0120$।
दोनों 1.2 प्रतिशत देते हैं।

\textbf{18.} $f = 0.75\times(-0.00019)/\bigl(40\times(-0.00147) + 0.75\times(-0.00019)\bigr)
= 0.0024$, जो Sr \qty{10.46}{ppm} और $^{87}\mathrm{Sr}/^{86}\mathrm{Sr} =
0.70361$ पूर्वानुमानित करता है, मापे गए 0.70290 के सामने। आशा प्रावार--अवसाद मिश्रण-वक्र पर नहीं है:
उसके स्रोत में एक अन्य घटक है।

\textbf{19.} पुरानी पुनर्चक्रित महासागरीय भूपर्पटी, HIMU प्रकार: उसके Nd समस्थानिकों को \omterm{def:b2:isotope-geochemistry:reservoirs}{क्षीणित प्रावार} से
कुछ निम्न Sm/Nd चाहिए, उसके Sr समस्थानिकों को निम्न Rb/Sr, और उसके
सीसे को (यहाँ मापा नहीं गया) ऊँचा U/Pb, सब एक से दो अरब वर्षों तक।

\textbf{20.} $e^{\lambda_{147}\cdot2\times10^9} - 1 = 0.013133$। अवसाद: $0.51141 -
0.110\times0.013133 = 0.509965$; CHUR: $0.512638 - 0.1966\times0.013133 = 0.510056$;
$\varepsilon_{\mathrm{Nd}}(2\,\mathrm{Ga}) = (0.509965/0.510056 - 1)\times10^4 = -1.8$।

\textbf{21.} $f = 0.110/0.1966 - 1 = -0.440$; क्षेपण के बाद से $\varepsilon$
$25.0\times(-0.440)\times2.0 = -22.0$ बदला, अतः $\varepsilon(2\,\mathrm{Ga}) \approx -24.0 + 22.0
= -2.0$, सटीक मान के 0.2 इकाई के भीतर।

\textbf{22.} $e^{\lambda_{87}\cdot2\times10^9} - 1 = 0.028338$;
$0.72000 - 0.50\times0.028338 = 0.70583$।

\textbf{23.} क्षेपण के समय अवसाद का $\varepsilon_{\mathrm{Nd}} = -1.8$ था, उस समय केवल
0.7 अरब वर्ष पुरानी भूपर्पटी के $-4.7$ के निकट, और पुरानी भूपर्पटी के
प्रबल ऋणात्मक मानों से बहुत ऊपर: यह युवा भूपर्पटी के अपरदन से मेल खाता है, शायद
प्रावार-व्युत्पन्न कुछ पदार्थ के साथ मिश्रित। उस समय उसका $^{87}\mathrm{Sr}/^{86}\mathrm{Sr}$
0.7058 भी एक युवा भूपर्पटी-चट्टान का है।

\textbf{24.} वेल का स्रोत, जो सबसे अधिक समृद्ध द्वीप है,
\textbf{\qty[text-series-to-math=true]{1.2}{\percent}} पुनर्चक्रित अवसाद धारण करता है।
\end{solution}
